\documentclass[12pt,a4paper]{article}

% ═══════════════════════════════════════════════════════
% Настройки для LuaLaTeX
% ═══════════════════════════════════════════════════════
\usepackage{fontspec}
\usepackage{polyglossia}
\setmainlanguage{russian}
\setotherlanguage{english}

% Шрифты с поддержкой кириллицы (должны быть установлены в системе)
\setmainfont{Times New Roman}
\setsansfont{Times New Roman}
\setmonofont{Times New Roman}

\usepackage{geometry}
\geometry{a4paper, left=2.5cm, right=2.5cm, top=2.5cm, bottom=2.5cm}

\usepackage{amsmath}
\usepackage{amssymb}
\usepackage{graphicx}
\usepackage{xcolor}
\usepackage{hyperref}
\hypersetup{
    colorlinks=true,
    linkcolor=blue!50!black,
    urlcolor=blue!70!black,
    citecolor=green!50!black
}
\usepackage{booktabs}
\usepackage{longtable}
\usepackage{array}
\usepackage{enumitem}
\usepackage{titlesec}
\usepackage{setspace}
\setstretch{1.15}

% Оформление заголовков
\titleformat{\section}{\Large\bfseries\color{blue!40!black}}{\thesection.}{0.6em}{}
\titleformat{\subsection}{\large\bfseries\color{blue!30!black}}{\thesubsection.}{0.5em}{}
\titleformat{\subsubsection}{\normalsize\bfseries}{\thesubsubsection.}{0.4em}{}

% Листинги кода через minted (требует -shell-escape)
\usepackage{minted}
\setminted{
    style=vs,
    fontsize=\small,
    frame=single,
    framesep=2mm,
    breaklines=true,
    tabsize=4,
    bgcolor=gray!5
}
\usemintedstyle{friendly}

% Пользовательские команды
\newcommand{\code}[1]{\texttt{#1}}
\newcommand{\term}[1]{\textit{#1}}

\title{\textbf{Архитектура научно-образовательной платформы}\\[0.5em]
\large Итоговый документ по проектированию}
\author{Проект <<analysis>>}
\date{\today}

\begin{document}

\maketitle

\begin{abstract}
\noindent
Документ обобщает архитектурные решения, принятые в ходе проектирования
научно-образовательной платформы для преподавателя физики. Платформа
предназначена для автоматизации создания справочных материалов, задач,
презентаций и последующей публикации в веб-пространстве. Изложение
учитывает переход от локального монолита к микросервисной архитектуре
и сохранение воспроизводимости результатов без обязательного
использования нейросетей.
\end{abstract}

\tableofcontents
\newpage

% ═══════════════════════════════════════════════════════
\section{Общая концепция}
% ═══════════════════════════════════════════════════════

Платформа решает три взаимосвязанные задачи:

\begin{enumerate}[label=\textbf{\arabic*.}]
    \item \textbf{Генерация справочных материалов} (PDF) из базы
          определений и формул.
    \item \textbf{Генерация задач} с автоматическим вычислением ответов
          на основе символьных вычислений.
    \item \textbf{Генерация презентаций} (reveal.js) для диктантов
          и лекций, включая интерактивные элементы.
\end{enumerate}

\noindent
В перспективе --- публикация тех же материалов через веб-интерфейс
с сохранением воспроизводимости и единого источника истины.

% ═══════════════════════════════════════════════════════
\section{Принципы проектирования}
% ═══════════════════════════════════════════════════════

\subsection{Разделение слоёв}

Проект строится по классической схеме <<данные --- домен --- представление>>:

\begin{center}
\begin{tabular}{@{}p{3.2cm}p{5.5cm}p{5.5cm}@{}}
\toprule
\textbf{Слой} & \textbf{Ответственность} & \textbf{Технологии} \\
\midrule
Хранилище & Данные о концепциях, формулах, задачах, величинах & SQLite (локально), PostgreSQL (веб) \\
Домен & Типизированные модели, бизнес-логика, генерация & dataclass/Pydantic, SymPy, Pint \\
Представление & Рендеринг PDF, reveal.js, JSON/HTML & Jinja2, LuaLaTeX, reveal.js, FastAPI \\
\bottomrule
\end{tabular}
\end{center}

\noindent
Такое разделение позволяет позже заменить SQLite на PostgreSQL,
а LuaLaTeX-рендеринг --- на веб-API, не переписывая бизнес-логику.

\subsection{Единый источник истины}

Все сгенерированные материалы (PDF, презентации, задачи) должны
строиться из одной базы данных. Это гарантирует, что:
\begin{itemize}
    \item изменения в определении автоматически попадают во все
          производные документы;
    \item нет дублирования текста в разных шаблонах;
    \item можно воспроизвести любой материал через \code{git checkout}
          и повторный запуск генератора.
\end{itemize}

\subsection{Воспроизводимость без обязательных нейросетей}

Нейросети допустимы как \textbf{вспомогательный инструмент на этапе
разметки данных} (первичное заполнение контекстов, формулировки),
но \textbf{не участвуют в рантайме}. Все вычисления, ответы и
описания генерируются детерминированно.

% ═══════════════════════════════════════════════════════
\section{Структура проекта}
% ═══════════════════════════════════════════════════════

\begin{verbatim}
analysis/
├── .venv/                     # виртуальное окружение (не коммитится)
├── data/
│   ├── physics.db             # SQLite (локально)
│   ├── raw/                   # исходные данные
│   └── processed/             # обработанные данные
├── docs/                      # документация (Sphinx / MyST)
├── help/                      # личные заметки и шпаргалки
├── media/
│   ├── figures/               # выходные графики
│   ├── tables/                # выходные таблицы
│   └── schemes/               # схемы
├── notebooks/                 # Jupyter для исследований
├── scripts/
│   ├── create_db.py
│   ├── import_csv.py
│   ├── generate_pdf.py
│   ├── generate_quiz.py
│   └── generate_lecture.py
├── src/my_project/
│   ├── __init__.py
│   ├── database.py            # подключение и миграции
│   ├── repositories.py        # доступ к данным
│   ├── domain/                # dataclass-модели
│   │   ├── content.py         # Formula, Concept, Figure
│   │   ├── quantities.py      # PhysicalQuantity, QuantityUsage
│   │   └── tasks.py           # Task, GeneratedVariant
│   ├── generation/            # NLG-генерация описаний
│   │   ├── grammar.py         # pymorphy3-обёртка
│   │   └── templates.py
│   ├── render/
│   │   ├── latex.py
│   │   ├── reveal.py
│   │   └── web.py             # будущий API
│   └── templates/
│       ├── physics.tex
│       └── lecture_reveal.j2
├── tests/
├── .gitignore
├── pyproject.toml
└── README.md
\end{verbatim}

% ═══════════════════════════════════════════════════════
\section{Модель данных}
% ═══════════════════════════════════════════════════════

\subsection{Основные таблицы}

\begin{center}
\begin{tabular}{@{}p{3.5cm}p{9.5cm}@{}}
\toprule
\textbf{Таблица} & \textbf{Назначение} \\
\midrule
\code{concepts} & Определения (текст, раздел, уровень, пример) \\
\code{formulas} & Формулы (LaTeX, раздел, уровень, пример) \\
\code{variables} & Физические величины (символ, имя, единица, род) \\
\code{formula\_variable\_context} & Квалификаторы переменных в конкретных формулах \\
\code{tasks} & Задачи (шаблон условия, SymPy-формула, спецификация параметров) \\
\code{task\_concepts}, \code{task\_formulas} & Связи задач с концепциями и формулами \\
\code{figures} & Рисунки и схемы (путь к изображению, тип) \\
\code{sources} & Библиография \\
\bottomrule
\end{tabular}
\end{center}

\subsection{Таблица контекста переменных}

Ключевая идея: квалификатор (<<точечного заряда в данной точке поля>>)
хранится не в самой величине, а в паре <<формула---величина>>:

\begin{minted}{sql}
CREATE TABLE formula_variable_context (
    formula_id   INTEGER NOT NULL,
    variable_id  INTEGER NOT NULL,
    qualifier    TEXT,          -- родительный падеж
    location     TEXT,          -- обстоятельство
    detail_level INTEGER DEFAULT 1,
    PRIMARY KEY (formula_id, variable_id),
    FOREIGN KEY (formula_id)  REFERENCES formulas(id),
    FOREIGN KEY (variable_id) REFERENCES variables(id)
);
\end{minted}

\noindent
Это позволяет описывать \(\Pi\) в разных формулах по-разному,
не дублируя данные.

\subsection{Таблица задач}

\begin{minted}{sql}
CREATE TABLE tasks (
    task_id            INTEGER PRIMARY KEY AUTOINCREMENT,
    source             TEXT,   -- задачник, номер
    section            TEXT,
    subsection         TEXT,
    level              TEXT,   -- "school:10-11", "university:1-2"
    condition_template TEXT,   -- шаблон с {m}, {v}, ...
    answer_formula     TEXT,   -- SymPy-выражение
    parameter_spec     TEXT,   -- JSON: диапазоны и единицы
    units              TEXT,   -- единицы ответа
    notes              TEXT
);

CREATE TABLE generated_variants (
    variant_id INTEGER PRIMARY KEY AUTOINCREMENT,
    task_id    INTEGER NOT NULL,
    seed       INTEGER NOT NULL,
    parameters TEXT,           -- JSON
    answer     TEXT,           -- точное значение
    FOREIGN KEY (task_id) REFERENCES tasks(task_id)
);
\end{minted}

\noindent
Варианты генерируются <<на лету>>, а в БД сохраняются только
зафиксированные варианты (например, для контрольных), с seed
для воспроизводимости.

% ═══════════════════════════════════════════════════════
\section{Доменная модель}
% ═══════════════════════════════════════════════════════

\subsection{Типы контента: ООП с наследованием}

Общие атрибуты (уровень, раздел, пример) вынесены в базовый класс,
наследники добавляют специфику.

\begin{minted}{python}
from dataclasses import dataclass, field
from typing import Optional, List

@dataclass
class ContentItem:
    item_id: int
    section: str
    subsection: str
    level: str
    example: Optional[str] = None
    notes: Optional[str] = None

@dataclass
class Formula(ContentItem):
    latex: str = ""
    variable_keys: List[str] = field(default_factory=list)

@dataclass
class Concept(ContentItem):
    definition: str = ""
    related_formula_ids: List[int] = field(default_factory=list)

@dataclass
class Figure(ContentItem):
    image_path: str = ""
    figure_type: str = "scheme"  # "scheme" | "plot"
    caption: str = ""
\end{minted}

\noindent
Обратите внимание: \code{Figure} хранит тип как поле, а не как
наследника --- это позволяет избежать <<ромбовидного наследования>>.

\subsection{Величины: композиция вместо наследования}

Наследование <<скорость \(\to\) скорость тела / скорость процесса>>
приводит к жёсткой иерархии: одна и та же величина может играть
разные роли. Поэтому используется композиция.

\begin{minted}{python}
@dataclass
class PhysicalQuantity:
    key: str
    symbol: str
    base_name: str
    unit: str
    gender: str                # "f", "m", "n"
    government: Optional[str]  # требуемый падеж дополнения

@dataclass
class QuantityUsage:
    quantity: PhysicalQuantity
    formula_id: int
    qualifier: Optional[str] = None
    location: Optional[str] = None
    detail_level: int = 1
\end{minted}

\noindent
Уровни детализации:
\begin{itemize}
    \item \textbf{1} --- только \code{base_name};
    \item \textbf{2} --- добавляется \code{qualifier};
    \item \textbf{3} --- добавляется \code{location}.
\end{itemize}

% ═══════════════════════════════════════════════════════
\section{Генерация текста (NLG)}
% ═══════════════════════════════════════════════════════

\subsection{Шаблоны с грамматическими метаданными}

Каждый слот несёт не только текст, но и грамматическую информацию:

\begin{minted}{python}
@dataclass
class Slot:
    text: str
    gender: str      # "f", "m", "n"
    number: str      # "sg", "pl"
    case: str        # "nom", "gen", "dat", ...
\end{minted}

\noindent
Согласование прилагательного с существительным, управление
падежом дополнения и склонение --- через \code{pymorphy3}.

\subsection{Пример генерации}

\begin{minted}{python}
import pymorphy3
morph = pymorphy3.MorphAnalyzer()

def inflect(phrase: str, case: str) -> str:
    words = phrase.split()
    result = []
    for w in words:
        p = morph.parse(w)[0]
        inflected = p.inflect({case})
        result.append(inflected.word if inflected else w)
    return " ".join(result)

def describe(usage: QuantityUsage) -> str:
    q = usage.quantity
    if usage.detail_level == 1 or not usage.qualifier:
        return f"{q.symbol} — {q.base_name}, {q.unit}"
    qualifier = inflect(usage.qualifier, "gen")
    parts = [q.base_name, qualifier]
    if usage.detail_level >= 3 and usage.location:
        parts.append(usage.location)
    return f"{q.symbol} — {' '.join(parts)}, {q.unit}"
\end{minted}

\noindent
Результат для \(\Pi\) в формуле потенциала:
\begin{quote}
\code{\textbackslash Pi — потенциальная энергия точечного заряда, Дж}
\end{quote}

\subsection{Ограничения подхода}

\begin{itemize}
    \item \code{pymorphy3} хорошо склоняет отдельные слова и простые
          словосочетания, но не строит сложные причастные обороты.
    \item Для фраз вида <<заряда, помещённого в данную точку поля>>
          рекомендуется хранить готовую формулировку в нужном падеже
          (в поле \code{qualifier}).
\end{itemize}

% ═══════════════════════════════════════════════════════
\section{Вычисления и единицы измерения}
% ═══════════════════════════════════════════════════════

\subsection{SymPy: символьные вычисления и ОДЗ}

Формулы ответов хранятся в виде SymPy-выражений. Это даёт
точные ответы и возможность автоматической проверки.

\begin{minted}{python}
import sympy as sp
from sympy.calculus.util import continuous_domain, singularities

x = sp.symbols('x', real=True)
expr = 1 / (x**2 - 1)

domain = continuous_domain(expr, x, sp.S.Reals)
sing = singularities(expr, x, sp.S.Reals)
\end{minted}

\noindent
Ограничения на параметры (например, \(m > 0\)) хранятся в
\code{parameter\_spec} и проверяются при генерации вариантов.

\subsection{Pint: единицы измерения}

Числовые расчёты с единицами выполняются через \code{Pint}:

\begin{minted}{python}
import pint
ureg = pint.UnitRegistry()

v = (50 * ureg.centimeter).to(ureg.meter) / (2 * ureg.second)
print(v.to(ureg.km / ureg.hour))
\end{minted}

\noindent
\textbf{SymPy} используется для символьного вывода (формулы, LaTeX),
\textbf{Pint} --- для проверки согласованности единиц и числовых
пересчётов.

% ═══════════════════════════════════════════════════════
\section{Генерация PDF (LuaLaTeX)}
% ═══════════════════════════════════════════════════════

\subsection{Почему LuaLaTeX}

Классический \code{pdflatex} требует дополнительных пакетов
(\code{inputenc}, \code{T2A}) и хуже работает с UTF-8 кириллицей
и системными шрифтами. \textbf{LuaLaTeX} нативно поддерживает UTF-8,
\code{fontspec} и \code{polyglossia}, что критично для русскоязычных
научных документов.

\subsection{Шаблон}

Шаблон \code{templates/physics.tex} использует:

\begin{minted}{latex}
\usepackage{fontspec}
\usepackage{polyglossia}
\setmainlanguage{russian}
\setmainfont{Liberation Serif}
\end{minted}

\noindent
Jinja2 подставляет данные из БД:

\begin{minted}{latex}
\begin{equation*}
{{ formula.formula_latex }}
\end{equation*}

\begin{itemize}
{% for usage in formula.usages %}
    \item[$\bullet$] {{ usage.description }}
{% endfor %}
\end{itemize}
\end{minted}

\subsection{Проверка компиляции}

Генератор \textbf{обязан проверять код возврата} \code{lualatex}
и не удалять временные файлы при ошибке:

\begin{minted}{python}
import shutil, subprocess
engine = shutil.which('lualatex') or shutil.which('xelatex')

for i in range(2):
    result = subprocess.run(
        [engine, '-interaction=nonstopmode', 'output.tex'],
        capture_output=True, text=True
    )
    if result.returncode != 0:
        print(result.stdout[-3000:])
        raise RuntimeError("LaTeX compilation failed")
\end{minted}

% ═══════════════════════════════════════════════════════
\section{Презентации (reveal.js)}
% ═══════════════════════════════════════════════════════

Презентации генерируются из того же источника данных через
Jinja2-шаблон. Для диктантов порядок случайный, для лекций ---
упорядоченный.

\subsection{Структура лекции}

\begin{minted}{sql}
CREATE TABLE lectures (
    lecture_id INTEGER PRIMARY KEY,
    title      TEXT NOT NULL,
    section    TEXT,
    level      TEXT
);

CREATE TABLE lecture_items (
    lecture_id   INTEGER NOT NULL,
    item_type    TEXT NOT NULL,  -- 'concept' | 'formula' | 'task'
    item_id      INTEGER NOT NULL,
    order_index  INTEGER NOT NULL,
    PRIMARY KEY (lecture_id, item_type, item_id),
    FOREIGN KEY (lecture_id) REFERENCES lectures(lecture_id)
);
\end{minted}

\subsection{Генерация слайдов}

\begin{minted}{python}
def build_lecture(lecture_id):
    items = get_lecture_items(lecture_id)
    slides = []
    for item in items:
        if item.type == 'concept':
            slides.append(render_concept_slide(item))
        elif item.type == 'formula':
            slides.append(render_formula_slide(item))
        elif item.type == 'task':
            slides.append(render_task_slide(item))
    return render_reveal_template(slides)
\end{minted}

% ═══════════════════════════════════════════════════════
\section{Микросервисная архитектура}
% ═══════════════════════════════════════════════════════

\subsection{Принципы}

Переход к микросервисам оправдан, когда появляется веб-интерфейс
и несколько сценариев использования. Разделение на сервисы
основано на \textbf{доменных границах}, а не на типах файлов.

\subsection{Сервисы}

\begin{center}
\begin{tabular}{@{}p{3.5cm}p{6cm}p{4cm}@{}}
\toprule
\textbf{Сервис} & \textbf{Ответственность} & \textbf{Стек} \\
\midrule
\code{content-service} & CRUD для концепций, формул, величин & FastAPI + PostgreSQL \\
\code{task-service} & Генерация задач, проверка ответов & FastAPI + SymPy + Pint \\
\code{render-service} & Генерация PDF и reveal.js & FastAPI + LuaLaTeX + Jinja2 \\
\code{web-frontend} & Веб-интерфейс для преподавателя и студента & React / Vue / Svelte \\
\code{auth-service} & Аутентификация и роли & FastAPI + OAuth2 \\
\bottomrule
\end{tabular}
\end{center}

\subsection{Взаимодействие}

\begin{itemize}
    \item \textbf{Синхронно} (REST) --- когда нужен немедленный ответ
          (например, <<сгенерируй PDF>>).
    \item \textbf{Асинхронно} (очередь, например, RabbitMQ или Kafka) ---
          для длительных операций (пакетная генерация, компиляция
          LaTeX, отчёты).
\end{itemize}

\subsection{Почему не монолит}

Монолит на FastAPI тоже возможен и проще на старте. Микросервисы
оправданы, когда:
\begin{itemize}
    \item есть разные требования к масштабированию (рендеринг
          PDF требует LuaLaTeX-окружения, задачи --- только SymPy);
    \item разные команды или пользователи работают с разными частями;
    \item нужно независимое развёртывание (например, обновить
          рендер, не трогая генерацию задач).
\end{itemize}

\noindent
\textbf{Рекомендация:} начать с \textbf{модульного монолита} на
FastAPI с чётким разделением слоёв. Выделять сервисы по мере
роста нагрузки. Это снижает риск преждевременной оптимизации.

% ═══════════════════════════════════════════════════════
\section{Веб-расширение}
% ═══════════════════════════════════════════════════════

\subsection{API}

Все операции, доступные локальным скриптам, должны быть
доступны через REST API. Это ключевой принцип: локальный
и веб-режим используют \textbf{один и тот же доменный слой}.

\begin{minted}{python}
# Пример будущего API
@app.get("/api/formulas")
def list_formulas(section: str = None, level: str = None):
    return repositories.get_formulas(section, level)

@app.post("/api/tasks/generate")
def generate_task(task_id: int, seed: int = None):
    return task_service.generate(task_id, seed)
\end{minted}

\subsection{Безопасность}

\begin{itemize}
    \item OAuth2 / JWT для аутентификации.
    \item Ролевая модель: \code{teacher}, \code{student}, \code{admin}.
    \item Ограничение доступа к генерации и редактированию
          контента.
\end{itemize}

\subsection{Хранение}

\begin{itemize}
    \item Локально --- \textbf{SQLite}.
    \item В вебе --- \textbf{PostgreSQL} (через SQLAlchemy + Alembic).
    \item Файлы (PDF, изображения) --- объектное хранилище (S3-совместимое,
          например, MinIO).
\end{itemize}

\subsection{Рендеринг в вебе}

LuaLaTeX-рендеринг в вебе --- тяжёлая операция. Варианты:
\begin{itemize}
    \item отдельный \textbf{render-service} с LuaLaTeX в контейнере;
    \item очередь задач (Celery + Redis) для пакетной генерации;
    \item кэширование готовых PDF по хэшу входных данных.
\end{itemize}

% ═══════════════════════════════════════════════════════
\section{Инструменты и зависимости}
% ═══════════════════════════════════════════════════════

\begin{center}
\begin{tabular}{@{}p{4cm}p{9.5cm}@{}}
\toprule
\textbf{Область} & \textbf{Инструменты} \\
\midrule
Хранение & SQLite (локально), PostgreSQL + SQLAlchemy + Alembic (веб) \\
Символьные вычисления & SymPy \\
Единицы измерения & Pint, sympy.physics.units \\
Морфология & pymorphy3 \\
Шаблонизация & Jinja2 \\
PDF & LuaLaTeX (fontspec, polyglossia), либо XeLaTeX \\
Презентации & reveal.js \\
Веб-фреймворк & FastAPI + Uvicorn \\
Очереди & Celery + Redis (при необходимости) \\
Тестирование & pytest, pytest-cov, mypy \\
Версионный контроль & Git (GitHub/GitLab) \\
Среда & venv, pyproject.toml, ruff, black \\
\bottomrule
\end{tabular}
\end{center}

\subsection{Почему именно LuaLaTeX}

Для русскоязычных научных документов \textbf{LuaLaTeX} предпочтительнее
\code{pdflatex}:
\begin{itemize}
    \item нативная поддержка UTF-8;
    \item \code{fontspec} и \code{polyglossia} для системных шрифтов;
    \item корректная работа с формулами и кириллицей без
          <<костылей>> с \code{T2A}.
\end{itemize}

\noindent
Установка: официальный установщик TUG с последующим использованием
\code{tlmgr install <package>} для добавления пакетов. Это
гибче, чем \code{apt install texlive-full}.

% ═══════════════════════════════════════════════════════
\section{Тестирование и качество}
% ═══════════════════════════════════════════════════════

\subsection{Уровни тестов}

\begin{enumerate}
    \item \textbf{Юнит-тесты} для чистых функций:
          \code{parse_intervals}, \code{inflect}, \code{describe},
          \code{calculate\_answer}.
    \item \textbf{Интеграционные тесты} для репозиториев на
          временной БД SQLite.
    \item \textbf{Снэпшот-тесты} для сгенерированного LaTeX
          (сравнение с эталоном).
    \item \textbf{Моки} для \code{subprocess.run} (LuaLaTeX)
          и \code{argparse}.
\end{enumerate}

\subsection{Статический анализ}

\begin{itemize}
    \item \code{mypy} --- проверка типов (PEP 484).
    \item \code{ruff} --- линтер и форматтер.
    \item \code{black} --- единый стиль (если не используете ruff format).
\end{itemize}

\subsection{CI/CD}

\begin{itemize}
    \item \code{pytest} + \code{mypy} + \code{ruff} на каждый push.
    \item Сборка Docker-образов для сервисов.
    \item Деплой через GitOps (например, ArgoCD) при переходе
          на Kubernetes.
\end{itemize}

% ═══════════════════════════════════════════════════════
\section{Дорожная карта}
% ═══════════════════════════════════════════════════════

\subsection{Этап 1: монолит (локально)}

\begin{itemize}
    \item Разделение слоёв (\code{repositories}, \code{domain},
          \code{render}).
    \item Добавление таблиц задач и контекстов переменных.
    \item Генерация PDF через LuaLaTeX с проверкой кода возврата.
    \item Генерация reveal.js для диктантов.
    \item Написание юнит-тестов для чистых функций.
\end{itemize}

\subsection{Этап 2: задачи с SymPy}

\begin{itemize}
    \item Генератор вариантов на SymPy с seed.
    \item Проверка ОДЗ и ограничений на параметры.
    \item Сохранение зафиксированных вариантов.
\end{itemize}

\subsection{Этап 3: модульный монолит на FastAPI}

\begin{itemize}
    \item REST API для контента и задач.
    \item Аутентификация и роли.
    \item Миграция на PostgreSQL.
    \item Веб-интерфейс (фронтенд).
\end{itemize}

\subsection{Этап 4: микросервисы}

\begin{itemize}
    \item Выделение \code{render-service} (тяжёлый LuaLaTeX).
    \item Выделение \code{task-service} (SymPy + Pint).
    \item Очередь задач для пакетной генерации.
\end{itemize}

\subsection{Этап 5: масштабирование}

\begin{itemize}
    \item Kubernetes, GitOps, мониторинг (Prometheus + Grafana).
    \item Кэширование PDF по хэшу.
    \item CDN для статики и изображений.
\end{itemize}

% ═══════════════════════════════════════════════════════
\section{Заключение}
% ═══════════════════════════════════════════════════════

Проект <<analysis>> сочетает сильные стороны локальной
воспроизводимости и современной веб-архитектуры. Ключевые решения:

\begin{enumerate}
    \item \textbf{Единый источник истины} --- база данных, из которой
          порождаются PDF, презентации и задачи.
    \item \textbf{Разделение слоёв} --- данные, домен, представление;
          позволяет заменить SQLite на PostgreSQL, а LuaLaTeX на API
          без переписывания логики.
    \item \textbf{Композиция вместо наследования} для величин;
          \textbf{наследование с оговорками} для типов контента.
    \item \textbf{Воспроизводимость без обязательных нейросетей}:
          SymPy, Pint, pymorphy3 и шаблоны с грамматическими
          метаданными.
    \item \textbf{Микросервисы по доменным границам}, а не по типам
          файлов; начать с модульного монолита на FastAPI.
    \item \textbf{LuaLaTeX} для русскоязычных PDF: нативная поддержка
          UTF-8, \code{fontspec}, \code{polyglossia}.
\end{enumerate}

\noindent
Такая архитектура позволяет сохранить локальную воспроизводимость
для научной работы, обеспечить готовность к веб-публикации и
избежать зависимости от нейросетей в критичных по корректности
участках.

\end{document}